% =============================================================================
%  WORSE IS BETTER  —  an editorial essay
%  Build:  lualatex worse-is-better.tex   (run twice for the TOC)
% =============================================================================
\documentclass[11pt]{article}

% ---- engine / fonts ---------------------------------------------------------
\usepackage{fontspec}
\usepackage[math-style=ISO]{unicode-math}

\setmainfont{EB Garamond}[
  Numbers = OldStyle, Ligatures = {TeX, Common},
]
\setsansfont{Fira Sans}[
  UprightFont = *, BoldFont = * SemiBold, Ligatures = TeX,
]
\newfontfamily\firalight{Fira Sans Light}
\newfontfamily\firabook{Fira Sans Book}
\setmonofont{Menlo}[Scale=0.86]
\setmathfont{Libertinus Math}

% ---- geometry ---------------------------------------------------------------
\usepackage[
  a4paper,
  inner=2.2cm, outer=4.9cm,
  top=2.4cm, bottom=2.5cm,
  marginparwidth=3.4cm, marginparsep=0.6cm, headsep=1.0cm,
]{geometry}

\usepackage{microtype}
\usepackage{graphicx}
\usepackage[table]{xcolor}
\usepackage{tikz}
\usetikzlibrary{calc,arrows.meta,decorations.pathmorphing,shadings,positioning}
\usepackage{enumitem}
\usepackage{booktabs}
\usepackage{array}
\usepackage{longtable}
\usepackage{lettrine}
\usepackage{fontawesome5}
\usepackage{ragged2e}
\usepackage[most,skins,breakable]{tcolorbox}
\usepackage{titlesec}
\usepackage{fancyhdr}
\usepackage{marginnote}
\usepackage[hidelinks]{hyperref}

% ---- palette (muted editorial) ---------------------------------------------
\definecolor{ink}{HTML}{2B2A28}       % warm near-black
\definecolor{paper}{HTML}{F7F3EA}     % cream
\definecolor{paperdeep}{HTML}{EFE9DA} % deeper cream (panels)
\definecolor{terra}{HTML}{A6432E}     % terracotta accent (New Jersey)
\definecolor{terrasoft}{HTML}{F0E0D8}
\definecolor{slate}{HTML}{3C5A6E}     % slate counter-accent (MIT)
\definecolor{slatesoft}{HTML}{E1E7EB}
\definecolor{brass}{HTML}{9C7A32}     % brass hairline
\definecolor{sage}{HTML}{4A6152}      % agentty green
\definecolor{sagesoft}{HTML}{E4EBE4}
\definecolor{faint}{HTML}{8B8578}     % muted caption grey
\definecolor{hair}{HTML}{CFC7B4}

\color{ink}
\pagecolor{paper}
\hypersetup{
  colorlinks=true, linkcolor=terra!85!black, urlcolor=slate, citecolor=terra,
  pdftitle={Worse is Better}, pdfauthor={after Richard P. Gabriel},
}

% ---- section styling: small-caps, numbered, hairline ------------------------
\titleformat{\section}
  {\normalfont}
  {}{0pt}
  {\bfseries\sffamily\large\color{terra}\thesection\enspace}
  [\vspace{-0.35em}{\color{hair}\titlerule[0.7pt]}]
\titleformat{\subsection}
  {\sffamily\bfseries\normalsize\color{ink!85}}
  {\thesubsection}{0.55em}{}
\titlespacing*{\section}{0pt}{1.7\baselineskip}{0.7\baselineskip}
\titlespacing*{\subsection}{0pt}{1.0\baselineskip}{0.3\baselineskip}

% wrap section title in sans small-caps look
\let\oldsection\section
\renewcommand{\section}[1]{\oldsection{\sffamily\bfseries\large #1}}

% ---- running heads (masthead style) -----------------------------------------
\pagestyle{fancy}
\fancyhf{}
\renewcommand{\headrulewidth}{0pt}
\fancyhead[L]{\firalight\footnotesize\color{faint}Worse is Better}
\fancyhead[R]{\firalight\footnotesize\color{faint}\thepage}
\fancyfoot[C]{\color{hair}\rule{2cm}{0.4pt}}

% ---- margin note ------------------------------------------------------------
\renewcommand*{\marginfont}{\firabook\footnotesize\color{faint}}
\newcommand{\mnote}[1]{\marginnote{\RaggedRight #1}}
\newcommand{\mnum}[2]{\marginnote{\RaggedRight
  {\sffamily\bfseries\color{terra}#1}\\[1pt]#2}}

% ---- drop cap ---------------------------------------------------------------
\newcommand{\dropcap}[2]{%
  \lettrine[lines=3, lhang=0.10, loversize=0.10, findent=3pt, nindent=5pt,
    slope=0pt, findent=2pt]{\color{terra}\sffamily\bfseries #1}{#2}}

% ---- pull quote (big, hanging, brass rule) ----------------------------------
\newcommand{\pull}[2]{%
  \par\vspace{0.5\baselineskip}
  \noindent\hspace*{-0.3cm}%
  \begin{minipage}{\dimexpr\linewidth+0.3cm}
    {\color{brass}\rule{2.4cm}{1.2pt}}\\[4pt]
    {\firalight\Large\itshape\color{terra!90!black}\setlength{\parindent}{0pt}#1}\\[5pt]
    {\sffamily\footnotesize\color{faint}#2}
  \end{minipage}\par\vspace{0.6\baselineskip}}

% ---- callouts ---------------------------------------------------------------
% generic side-tab box
\newtcolorbox{sidebox}[2][]{%
  enhanced, breakable, sharp corners,
  colback=#2!8, colframe=#2!8, boxrule=0pt,
  borderline west={3pt}{0pt}{#2},
  left=13pt, right=11pt, top=9pt, bottom=9pt,
  fontupper=\normalsize,
  overlay unbroken and first={%
    \node[anchor=north west, font=\sffamily\bfseries\footnotesize\color{#2!80!black}]
      at ([xshift=13pt,yshift=-5pt]frame.north west) {#1};},
  before upper={\vspace{0.75\baselineskip}},
}
\newenvironment{keybox}[1][The idea]{\begin{sidebox}[#1]{terra}}{\end{sidebox}}
\newenvironment{counterbox}[1][Counterpoint]{\begin{sidebox}[\faBalanceScale\ \ #1]{slate}}{\end{sidebox}}
\newenvironment{agentbox}[1][In agentty]{\begin{sidebox}[\faTerminal\ \ #1]{sage}}{\end{sidebox}}
\newenvironment{mythbox}[1][But wait]{\begin{sidebox}[\faQuestionCircle\ \ #1]{brass}}{\end{sidebox}}

% "case file" framed box for the marquee story
\newtcolorbox{casefile}[1]{%
  enhanced, breakable,
  colback=paperdeep, colframe=ink!55, boxrule=0.8pt,
  arc=0pt, outer arc=0pt,
  left=14pt, right=14pt, top=30pt, bottom=12pt,
  overlay unbroken and first={%
    \fill[ink!55] (frame.north west) rectangle
      ([yshift=-20pt]frame.north east);
    \node[anchor=west, font=\sffamily\bfseries\small\color{paper}]
      at ([xshift=12pt,yshift=-10pt]frame.north west)
      {\faFolderOpen\ \ CASE FILE\, ---\, #1};},
}

% ---- figure caption helper --------------------------------------------------
\newcommand{\figcap}[2]{\par\vspace{2pt}\noindent
  {\sffamily\footnotesize\color{terra}\textbf{Fig.\ #1}}\,%
  {\firabook\footnotesize\color{faint}\enspace#2}\par}

% ---- misc -------------------------------------------------------------------
\setlength{\parindent}{1.15em}
\setlength{\parskip}{0pt}
\setlist{itemsep=2.5pt, topsep=4pt, parsep=0pt}
\renewcommand{\arraystretch}{1.3}
\newcommand{\NJ}{\textcolor{terra}{\sffamily\bfseries NEW JERSEY}}
\newcommand{\MITs}{\textcolor{slate}{\sffamily\bfseries MIT}}
\newcommand{\wib}{\textit{worse\,is\,better}}

% small dinkus separator
\newcommand{\dinkus}{\par\vspace{4pt}
  \centerline{\color{brass}\small\quad$\ast\ \ \ast\ \ \ast$\quad}\vspace{4pt}\par}

% =============================================================================
\begin{document}

% ========================================================= COVER =============
\begin{titlepage}
\hypersetup{pageanchor=false}
\thispagestyle{empty}
\begin{tikzpicture}[remember picture, overlay]

  % full-bleed cream
  \fill[paper] (current page.south west) rectangle (current page.north east);

  % ---- masthead: thin double rule + kicker ----
  \coordinate (TL) at ([yshift=-2.6cm,xshift=2.2cm]current page.north west);
  \coordinate (TR) at ([yshift=-2.6cm,xshift=-2.2cm]current page.north east);
  \draw[ink, line width=1.1pt] (TL) -- (TR);
  \draw[ink, line width=0.5pt] ($(TL)+(0,-0.14)$) -- ($(TR)+(0,-0.14)$);
  \node[anchor=south, font=\sffamily\footnotesize, text=faint]
    at ($(TL)!0.5!(TR)+(0,0.08)$)
    {\textls[300]{AN ESSAY ON HOW SOFTWARE ACTUALLY WINS}};

  % ---- big title ----
  \node[anchor=north, align=center] at ([yshift=-3.5cm]current page.north)
    {\firalight\fontsize{58}{58}\selectfont\color{ink} Worse\\[-6pt]};
  \node[anchor=north, align=center] at ([yshift=-6.0cm]current page.north)
    {\itshape\fontsize{30}{30}\selectfont\color{brass} is};
  \node[anchor=north, align=center] at ([yshift=-7.2cm]current page.north)
    {\sffamily\bfseries\fontsize{58}{58}\selectfont\color{terra} Better};

  % ---- deck / standfirst ----
  \node[anchor=north, align=center, text width=12cm]
    at ([yshift=-10.3cm]current page.north)
    {\firabook\large\color{ink!85}
     Why the simple, good-enough thing that ships \emph{today}\\
     so often beats the perfect thing that ships \emph{tomorrow}\,---\\
     and what a Lisp hacker's 1989 lament still teaches us.};

  % ---- the "vs" diptych, refined ----
  \begin{scope}[shift={($(current page.center)+(0,-1.6cm)$)}]
    % MIT plate
    \fill[slate] (-3.55,-1.35) rectangle (-0.25,1.35);
    \draw[paper, line width=0.5pt] (-3.4,-1.2) rectangle (-0.4,1.2);
    \node[paper, font=\sffamily\bfseries\LARGE] at (-1.9,0.52) {MIT};
    \node[slatesoft, font=\firabook\small, align=center]
      at (-1.9,-0.45) {\itshape ``the Right Thing''};
    % NJ plate
    \fill[terra] (0.25,-1.35) rectangle (3.55,1.35);
    \draw[paper, line width=0.5pt] (0.4,-1.2) rectangle (3.4,1.2);
    \node[paper, font=\sffamily\bfseries\LARGE] at (1.9,0.52) {N\,J};
    \node[terrasoft, font=\firabook\small, align=center]
      at (1.9,-0.45) {\itshape ``Worse is Better''};
    % vs monogram
    \node[circle, fill=paper, draw=ink, line width=0.8pt, minimum size=0.95cm,
          font=\itshape\normalsize\color{ink}] at (0,0) {vs};
  \end{scope}

  % ---- footer colophon ----
  \coordinate (BL) at ([yshift=2.5cm,xshift=2.2cm]current page.south west);
  \coordinate (BR) at ([yshift=2.5cm,xshift=-2.2cm]current page.south east);
  \draw[ink, line width=0.5pt] (BL) -- (BR);
  \node[anchor=north, align=center] at ($(BL)!0.5!(BR)+(0,-0.22)$)
    {\sffamily\footnotesize\color{ink}
     after \textbf{Richard P.\ Gabriel} $\cdot$
     \emph{Lisp: Good News, Bad News, How to Win Big} (1989)\\[3pt]
     {\color{faint}\firabook a reading of the essay, its critics,
     and a field guide for the \texttt{agentty} coding assistant}};
\end{tikzpicture}
\end{titlepage}

% ========================================================= FRONT MATTER ======
\clearpage
\hypersetup{pageanchor=true}
\pagenumbering{roman}
\thispagestyle{empty}

\vspace*{0.3cm}
\noindent{\color{hair}\rule{\linewidth}{0.7pt}}\\[-2pt]
\noindent{\sffamily\bfseries\Large\color{terra}Contents}\hfill
{\firalight\small\color{faint}nine short chapters}\\[-4pt]
\noindent{\color{hair}\rule{\linewidth}{0.7pt}}

\vspace{0.5cm}
\begin{minipage}{0.62\linewidth}
{\hypersetup{linkcolor=ink!85}\tableofcontents}
\end{minipage}\hfill
\begin{minipage}{0.33\linewidth}
\raggedright\firabook\small\color{ink!80}
\textbf{\sffamily\color{terra}In one sentence.}\\[3pt]
A simple, good-enough thing that ships \emph{today} usually beats a beautiful,
perfect thing that ships \emph{next year}\,---\,because by the time the perfect
thing arrives, everyone already uses (and is locked into) the simple one.
\end{minipage}

\clearpage

% =============================================================================
\pagenumbering{arabic}
\section{The idea, explained like you're new to it}
% =============================================================================

\dropcap{I}{magine} two restaurants opening on the same street. \textbf{A}
opens next week: small menu, burgers and fries, nothing fancy\,---\,but the
doors are \emph{open} and people can eat. \textbf{B} will be \emph{perfect}: a
Michelin chef, a wine cellar, hand-made everything. It opens \emph{in two
years}. Maybe.

\mnote{The whole theory fits inside a story about lunch.}
By the time B finally opens, A already has regulars, a loyalty card, a delivery
deal, and a reputation. The neighborhood \emph{eats at A}. B may be objectively
better\,---\,and it still loses, or claws for scraps. That, in a nutshell, is
\wib.

The phrase was coined by \textbf{Richard P.\ Gabriel}, a Lisp programmer, in a
1989 essay. He was chasing a puzzle that nagged at him: why do \emph{technically
inferior} pieces of software so often beat \emph{technically superior} ones? His
answer\,---\,because ``worse'' designs are simpler, so they ship sooner, spread
faster, and get entrenched before the ``right'' design is even finished.

\begin{keybox}[The core mechanism]
A simpler design is \emph{cheaper to build}, \emph{easier to move to new
machines}, and \emph{easier to change}. So it reaches users first. Once they
adopt it, three barriers lock them in: \textbf{habit} (they learned its quirks),
\textbf{compatibility} (their other tools now depend on it), and
\textbf{momentum} (everyone else uses it too). The ``better'' design arrives late
to a party where every guest already has a partner.
\end{keybox}

Gabriel pointed to real history where ``technically worse'' won: \textbf{C} beat
\textbf{Lisp}; \textbf{Unix} beat Lisp machines and \textbf{VMS}; \textbf{x86}
beat the cleaner \textbf{RISC} chips. He even called Unix and C ``the ultimate
computer viruses''\,---\,not an insult, but a compliment about how they
\emph{spread}.

\pull{As long as the initial program is good enough, it will spread\ldots\ long
before a program developed using the MIT approach has a chance to be
deployed.}{The argument in one breath}

% =============================================================================
\section{Where it came from --- a story that spread the way it describes}
% =============================================================================

Gabriel formed the idea in 1989 inside a longer essay,
\emph{``Lisp: Good News, Bad News, How to Win Big.''} One section was titled
\textbf{``The Rise of `Worse is Better'.''}

\mnote{Fittingly, the essay itself spread the \wib\ way\,---\,informally, virally,
good-enough, and fast.}
It could have stayed obscure. But in 1991 the programmer \textbf{Jamie Zawinski}
found it in Gabriel's files at Lucid Inc.\ and emailed it to friends. From there
it travelled by word of mouth: Zawinski's friends at Carnegie Mellon forwarded
it to friends at Bell Labs, ``who sent them to their friends everywhere.'' It
was reprinted as an appendix in the 1994 book \emph{The UNIX-HATERS Handbook},
and became famous enough to be credited with popularizing the ``East Coast vs.\
West Coast'' split among developers\,---\,a framing, we'll see, that Gabriel
never actually used.

% =============================================================================
\section{Two philosophies, face to face}
% =============================================================================

Gabriel described two schools of design, each with several names\,---\,which is
why the same idea appears under different labels.

\vspace{2pt}
\begin{center}
\begin{tikzpicture}[font=\sffamily]
  \node[fill=slatesoft, draw=slate, line width=0.9pt,
        text width=6.1cm, align=center, inner sep=9pt] (mit) at (-3.55,0)
    {\bfseries\color{slate}The ``MIT'' approach\\[2pt]
     \normalfont\firabook\footnotesize\color{ink}
     a.k.a.\ MIT/Stanford $\cdot$ ``The Right Thing''\\
     \emph{Common Lisp, Scheme}};
  \node[fill=terrasoft, draw=terra, line width=0.9pt,
        text width=6.1cm, align=center, inner sep=9pt] (nj) at (3.55,0)
    {\bfseries\color{terra}The ``New Jersey'' style\\[2pt]
     \normalfont\firabook\footnotesize\color{ink}
     a.k.a.\ Berkeley $\cdot$ ``Worse is Better''\\
     \emph{Unix \& C, out of Bell Labs, NJ}};
  \node[circle, fill=paper, draw=ink, line width=0.7pt, minimum size=0.8cm,
        font=\itshape\small] at (0,0) {vs};
\end{tikzpicture}
\end{center}
\figcap{1}{The two camps. Both prize the same four qualities; they rank them
differently, and the ranking is the whole ball game.}

\vspace{4pt}
Both schools care about \textbf{Simplicity}, \textbf{Correctness},
\textbf{Consistency}, and \textbf{Completeness}\,---\,in that descending order of
importance. They just disagree on what to sacrifice.

\vspace{4pt}
{\small
\begin{longtable}{@{}>{\sffamily\bfseries}p{2.0cm} >{\RaggedRight}p{5.4cm} >{\RaggedRight}p{5.4cm}@{}}
\toprule
& \cellcolor{slatesoft}\sffamily\bfseries\color{slate}MIT / The Right Thing
& \cellcolor{terrasoft}\sffamily\bfseries\color{terra}New Jersey \\
\midrule
\endhead
Simplicity &
Simple in both\,---\,but if forced, keep the \emph{interface} simple, even at the
cost of a hairier implementation. &
Simple in both\,---\,but if forced, keep the \emph{implementation} simple.
Simplicity is the single most important thing. \\[2pt]
Correctness &
``Incorrectness is simply not allowed.'' Non-negotiable. &
Correct in all observable ways, yet ``slightly better to be simple than
correct.'' \\[2pt]
Consistency &
Must be consistent; as important as correctness. May sacrifice a little
simplicity to keep it. &
Must not be \emph{overly} inconsistent. May trade consistency for simplicity;
drop rare cases rather than complicate the code. \\[2pt]
Completeness &
All reasonably expected cases \emph{must} be covered; simplicity may not
``overly reduce'' completeness. &
Cover what's practical, but completeness yields to \emph{any} other quality, and
\emph{must} yield when it threatens implementation simplicity. \\
\bottomrule
\end{longtable}}
\figcap{2}{Gabriel's own four-by-two table, reworded. Read the two Correctness
cells back to back: that single row is the whole disagreement.}

\begin{keybox}[The difference in one line]
For the \MITs{} mind, being \emph{wrong} is unacceptable\,---\,correctness first.
For the \NJ{} mind, being \emph{complicated} is unacceptable\,---\,simplicity
first, and a little wrongness at the edges is a fair price to ship.
\end{keybox}

% =============================================================================
\section{The case file: ``PC loser-ing'' --- the story at the heart of the essay}
% =============================================================================

Gabriel doesn't argue in the abstract. He tells a \emph{true story} about two
engineers arguing over one gnarly operating-system problem\,---\,and it's the
best single illustration of the whole idea. Here it is, unpacked.

\begin{casefile}{the interrupted system call}
\small
\textbf{The setup.} A program asks the operating system to do something slow
(say, read from disk into a buffer). Halfway through, an \textbf{interrupt}
fires\,---\,something else needs attention \emph{right now}. The system must
decide: what happens to the half-finished operation?

\smallskip
\textbf{\color{slate}The MIT answer (``the Right Thing'').} Quietly rewind. Back
the program's counter up to \emph{before} it made the call, save every scrap of
state, handle the interrupt, then transparently re-enter the call. The program
never even knows it was interrupted. Beautiful\,---\,and a nightmare to implement,
because the kernel must be able to unwind partial state for \emph{every}
long-running call. Gabriel's nickname for the mechanism: \textbf{PC loser-ing},
because the program counter (``PC'') gets coerced into ``loser mode''
(``loser'' being MIT's affectionate word for ``user'').

\smallskip
\textbf{\color{terra}The New Jersey answer (``Worse is Better'').} Don't rewind.
Just \emph{give up}, finish the call, and hand back an error code that means
``I got interrupted\,---\,you deal with it.'' A correct program must now check for
that code and \emph{retry the call itself}. Simpler kernel; more work pushed onto
every program that calls it.
\end{casefile}

\mnote{\faQuoteLeft\ ``Sometimes it takes a tough man to make a tender chicken,''
the MIT guy mutters at the end\,---\,and, Gabriel admits, ``the New Jersey guy
didn't understand. (I'm not sure I do either.)''}
The MIT engineer hated the New Jersey answer: the \emph{interface} is now
complex (every caller must loop-and-retry) even though the \emph{implementation}
is simple. The New Jersey engineer shrugged: that's the right trade\,---\,
implementation simplicity beats interface simplicity.

\subsection{The twist: what that error code became}
If you've ever written Unix code, you've met the New Jersey answer in person. It
is \textbf{\texttt{EINTR}} --- ``interrupted system call'' --- and to this day
robust programs wrap their slow calls in a retry loop:

\begin{center}\small\ttfamily
\begin{minipage}{9cm}
again:\\
\hspace*{2em}n = read(fd, buf, size);\\
\hspace*{2em}if (n < 0 \&\& errno == EINTR) goto again;
\end{minipage}
\end{center}

\mnote{\faLightbulb\ The ``worse'' system \emph{became} the right thing\,---\,
exactly as Gabriel predicted in step three. But the wart never fully healed.}
Here's the beautiful part, dug up by later engineers: Berkeley Unix added
\emph{automatic} retry (the \texttt{SA\_RESTART} flag) back in \textbf{1983}\,---\,
five years \emph{before} the essay was written. In other words, the ``worse''
system quietly grew ``the Right Thing'' on its own, in response to users tired
of writing that loop. And yet the scar remains: to this day, a careful Unix
programmer still has to think about \texttt{EINTR}. \emph{That} is \wib\ in one
real API: shipped simple, spread everywhere, improved later\,---\,but never
quite perfectly.

% =============================================================================
\section{Why the ``worse'' one wins --- the three-step trap}
% =============================================================================

Gabriel's case for \emph{why} it works is a chain of three steps. Slow down
here; the whole idea lives on this staircase.

\vspace{4pt}
\begin{center}
\begin{tikzpicture}[font=\sffamily\small, scale=0.94, transform shape]
  \def\bw{4.35}\def\bh{2.7}\def\dx{0.3}\def\dy{0.6}
  \foreach \i/\c/\t/\d in {%
    0/{terra!45}/{Ship first}/{Simple $\Rightarrow$ cheap, portable, fast to
       adapt\,---\,out spreading \emph{before} the ``right thing'' is done.},
    1/{terra!68}/{Lower the bar}/{Users adapt to its rough edges and stop
       expecting the ``right thing.'' It ``conditions users to expect less.''},
    2/{terra!90}/{Grow to almost-right}/{Now ubiquitous, it attracts the talent
       and pressure that grow it into \emph{almost} the right thing.}} {
    \pgfmathsetmacro{\x}{\i*(\bw+\dx) - (\bw+\dx)}
    \pgfmathsetmacro{\y}{\i*\dy}
    \pgfmathsetmacro{\tw}{\bw-0.55}
    \node[fill=\c, text=white, sharp corners,
          minimum width=\bw cm, minimum height=\bh cm,
          text width=\tw cm, align=center, inner sep=6pt] (s\i) at (\x,\y)
      {\bfseries\Large \the\numexpr\i+1\relax\\[2pt]
       \bfseries\normalsize \t\\[4pt]\normalfont\footnotesize \d};
  }
  \draw[-{Latex[length=2.6mm]}, ink!70, line width=1.1pt] (s0.east) -- (s1.west);
  \draw[-{Latex[length=2.6mm]}, ink!70, line width=1.1pt] (s1.east) -- (s2.west);
\end{tikzpicture}
\end{center}
\figcap{3}{The staircase. Each step is also a rung the ``right thing'' never gets
to climb, because it's still in the lab when step one happens.}

\mnote{\faLightbulb\ Popularity \emph{recruits} the people who make a tool
better. That is the quiet engine of step three.}
Gabriel's sharpest line lives in step three: \emph{even though Lisp compilers in
1987 were about as good as C compilers, there were far more experts trying to
improve C compilers than Lisp compilers}. Once a tool is everywhere, the crowd
that improves it is everywhere too.

\subsection{The ``50\%--80\% solution''}
Gabriel is precise about how good ``good enough'' must be: early Unix and C, he
says, ``provide about 50\%--80\% of what you want.'' That's the sweet spot\,---\,
enough to be genuinely useful, little enough to stay simple and portable. And
because half of all computers are always below the median in power, a lean 50\%
tool that \emph{runs everywhere} beats a lavish 100\% tool that needs a
supercomputer. He also names the failure modes of The Right Thing:

\begin{itemize}[leftmargin=1.4em]
  \item \textbf{The big complex system.} Aims for 100\%; ``the last 20\% takes
        80\% of the effort''; ships late, huge, and only on top-end hardware.
        \emph{(His example: Common Lisp.)}
  \item \textbf{The diamond-like jewel.} Stays tiny and elegant at every step\,---\,
        but making it run fast is ``beyond the capabilities of most
        implementors.'' \emph{(His example: Scheme.)}
\end{itemize}

He tied all this to \textbf{Richard Stallman}, to the philosophy of \textbf{Unix},
and to the \textbf{open-source movement}\,---\,all of which fed straight into
\textbf{Linux}. The pattern is not a one-off; it's a recurring shape.

% =============================================================================
\section{The counter-argument: ``Worse is Worse''}
% =============================================================================

Here is what most retellings leave out: serious people think Gabriel was
\emph{wrong}\,---\,and the sharpest of them is \textbf{Jim Waldo}, a Distinguished
Engineer at Sun, whose 2003 essay \emph{``Worse is Worse''} is a direct rebuttal.
His case is worth taking seriously.

\begin{counterbox}[Waldo's thesis]
``\emph{Worse is better} is a much catchier slogan than \emph{better depends on
your goodness metric}.'' The dominant technologies weren't \emph{worse}\,---\,they
were \emph{better on the axes customers actually cared about}. We just measured
``better'' with the wrong ruler.
\end{counterbox}

Run his argument through the famous examples:

\begin{itemize}[leftmargin=1.4em]
  \item \textbf{C over Lisp.} C ``produced faster code, was easier to master,
        was easier to use in groups, and ran well on less expensive hardware.''
        Those \emph{are} dimensions of better. On them, C winning was
        \emph{better is better}, not worse is better.
  \item \textbf{VHS over Betamax.} The classic ``inferior format won'' story\,---\,
        but VHS recorded \emph{longer} tapes on \emph{cheaper} machines from
        \emph{more} suppliers. Better on the axes buyers valued (length, cost),
        worse only on the one they didn't (picture quality).
  \item \textbf{Unix / DOS over Multics / VMS.} Unix ``had the twin advantages of
        actually existing, and running on a wide variety of cheap hardware.''
        \emph{Existence on cheap hardware} was itself the metric of goodness that
        won.
\end{itemize}

\mnote{Waldo also corrects a myth: Gabriel's essay \emph{never} says ``East
Coast vs.\ West Coast''\,---\,it says MIT/Stanford vs.\ New Jersey. The
coasts were added by retellers.}
Waldo's deeper worry is cultural. The slogan, he argues, ``has become a catch
phrase for those who\ldots\ never understood [the article] in the first
place''\,---\,used to ``justify shoddy design'' and to whisper that ``our
customers are too stupid to\ldots\ deserve high quality products.'' His verdict:
reject that reading and ``start putting out software that really is better\,---\,
on the dimension of goodness that our customers have.''

\pull{Better is a complicated notion, and can depend on a variety of different
metrics. What we geeks think of as better may not be what our customers think is
better.}{Jim Waldo, ``Worse is Worse'' (2003)}

And here's the twist that makes the whole debate honest: \textbf{Gabriel agreed
enough to argue against himself.} In 2000 he wrote \emph{``Worse Is Better Is
Worse''} under the pseudonym \emph{Nickieben Bourbaki}, and separately \emph{``Is
Worse Really Better?''}, applying the idea to C++'s triumph over more elegant
object-oriented languages. He has publicly called himself of two minds. So \wib\
is best read not as a \emph{commandment} but as a \emph{description of a
force}\,---\,one even its author distrusts.

% =============================================================================
\section{The family of related ideas}
% =============================================================================

The same wisdom, seen from other angles:

\begin{description}[leftmargin=1.4em, style=nextline, font=\sffamily\bfseries\color{ink!85}]
  \item[Gresham's law] ``Bad money drives out good''\,---\,cheap-and-common crowds
        out rare-and-fine.
  \item[If it ain't broke, don't fix it] Don't gold-plate what already works.
  \item[KISS \& Less is more] Simplicity as a first-class goal; fewer features
        done well beat many done poorly.
  \item[Minimum viable product] Ship the smallest thing that delivers value, then
        learn from real users\,---\,\wib\ as a startup method.
  \item[Neats vs.\ scruffies] In AI: clean theory (MIT) vs.\ things that work now
        (New Jersey).
  \item[Perfect is the enemy of good] Wait for perfect and you ship nothing.
  \item[Satisficing] Herbert Simon's word for accepting ``good enough'' instead
        of hunting the optimum\,---\,often the rational move under time limits.
  \item[Wabi-sabi] The Japanese aesthetic that finds beauty in the imperfect and
        incomplete\,---\,and Christopher Alexander's ``piecemeal growth,'' which
        Gabriel himself cites.
\end{description}

\subsection{``Wait --- isn't this the \emph{opposite} of KISS?''}
A fair objection. KISS says \emph{simple is good}; the title says \emph{worse} is
better. The knot unties once you see that \textbf{``worse'' is sarcasm.} Gabriel
isn't confessing low quality\,---\,he's needling the purists: \emph{the very thing
you call worse is the thing that wins.} Decode the two words and they stop
fighting:

\begin{itemize}[leftmargin=1.4em]
  \item \textbf{``Better''} $=$ survives, spreads, wins in the real world.
  \item \textbf{``Worse''} $=$ scores lower on the MIT scorecard\,---\,less
        \emph{complete}, not perfectly correct at every edge\,---\,\emph{not}
        lower quality overall.
\end{itemize}

Read that way, KISS and \wib\ aren't opposites\,---\,they're the \emph{same move},
\textbf{simplicity beats completeness}, differing only in emphasis:

\vspace{2pt}
{\small
\begin{longtable}{@{}>{\RaggedRight}p{2.7cm} >{\RaggedRight}p{5.0cm} >{\RaggedRight}p{5.2cm}@{}}
\toprule
& \sffamily\bfseries KISS & \sffamily\bfseries\color{terra}Worse is Better \\
\midrule
\endhead
The advice &
``Keep it simple.'' &
``Keep it simple \emph{even when that means leaving cases out and being slightly
wrong at the edges}.'' \\[2pt]
The tone &
A design \emph{virtue}\,---\,simple is nicer. &
A market \emph{prediction}\,---\,the simple one out-competes the complete one. \\[2pt]
Argues against &
Over-engineering. &
The MIT belief that you must be complete and perfectly correct \emph{before}
shipping. \\
\bottomrule
\end{longtable}}

\begin{mythbox}[The one honest difference]
KISS is purely about simplicity and never tells you to accept being \emph{wrong}.
\wib\ goes one notch further: a little incorrectness \emph{at the rare edges} is
an acceptable price for staying simple. But even here the first version must be
\textbf{good enough to work}. So it's KISS pushed one notch toward \emph{ship
it}\,---\,never an excuse for broken.
\end{mythbox}

% =============================================================================
\section{A field guide for \texorpdfstring{\texttt{agentty}}{agentty}}
% =============================================================================

\texttt{agentty} is a terminal coding assistant\,---\,and it is the perfect
crucible for this whole debate, because it has a \textbf{split personality}. Read
its own docs and the tension is right there on the page.

\begin{keybox}[The split personality]
The \emph{codebase} of \texttt{agentty} is a card-carrying \MITs{} ``Right
Thing'' project. But the \emph{behavior} it ships\,---\,how the binary acts, and
how the agent inside it should edit your code\,---\,is pure \NJ. Knowing which
hat to wear \emph{when} is the entire art of working on it.
\end{keybox}

\subsection{The codebase is ``The Right Thing'' --- on purpose}
Agentty's own \texttt{docs/DESIGN.md} opens with a line no New Jersey hacker
would write:

\pull{If a change makes the codebase look inconsistent\,---\,if it requires a
reader to ask ``wait, is this that style or this style?''\,---\,it doesn't go in.
There is one style.}{\texttt{docs/DESIGN.md}, ``the one idea''}

\mnote{Look at Gabriel's MIT column again: \emph{consistency as important as
correctness; incorrectness simply not allowed}. That \emph{is} agentty's design
contract.}
That is Gabriel's MIT scorecard almost verbatim. Concretely, agentty:

\begin{sloppypar}
\begin{itemize}[leftmargin=1.4em]
  \item makes illegal states \emph{unrepresentable}\,---\,every state machine is a
        \texttt{std::variant}, every fallible call an
        \texttt{std::expected<T,E>} (\texttt{DESIGN.md}'s seven rules);
  \item pins invariants at \emph{compile time}\,---\,151+ \texttt{static\_assert}
        / \texttt{consteval} proofs (\texttt{tool/policy.hpp} proves the
        permission matrix cell by cell; \texttt{provider/error\_class.hpp}
        proves every retry branch). ``The build is the test runner'': green build
        $\Rightarrow$ the invariants hold. \emph{Incorrectness is not allowed}
        because it \emph{will not compile};
  \item funnels all the messy dynamism (JSON, network, terminal IO) through
        \emph{one thin boundary}, past which everything is typed.
\end{itemize}
\end{sloppypar}

This is the opposite of ``ship it and check \texttt{EINTR} later.'' And it's the
\emph{right} call here, by our own decision guide: a coding agent that edits your
files and runs commands is exactly the high-stakes, irreversible domain where
Gabriel and Waldo both say \emph{do the Right Thing up front}.

\subsection{The behavior is ``Worse is Better'' --- also on purpose}
Now flip to what the binary \emph{does}. Every headline feature is a New Jersey
move:

\vspace{2pt}
{\small
\begin{longtable}{@{}>{\RaggedRight}p{5.4cm} >{\RaggedRight}p{7.8cm}@{}}
\toprule
\sffamily\bfseries agentty behavior & \sffamily\bfseries the \wib\ move it embodies \\
\midrule
\endhead
One 16.7\,MB native binary, {\raise0.1ex\hbox{$\sim$}}3\,ms cold start, zero runtime deps &
The lean 50--80\% tool that \emph{runs everywhere}, versus the lavish
Node/Python stack that needs its whole ecosystem installed first. \\
Retrieval sends only relevant slices ({\raise0.1ex\hbox{$\sim$}}80\% less context) &
``Good enough'' beats ``complete'': don't dump the whole repo, ship the slices
that matter and iterate. Falls back to keyword search when embeddings are
absent\,---\,\emph{always works, sometimes worse}. \\
Smart Mode scales effort to each turn &
A one-line rename and a cross-file refactor ``don't deserve the same effort.''
Spend the minimum that works; escalate only when a cheap attempt falls short. \\
\texttt{edit} preferred over \texttt{write} for changes &
The smallest diff that solves the problem\,---\,streams less, reviews easier,
reverts cleanly. \\
\bottomrule
\end{longtable}}
\figcap{4}{Agentty's product decisions, read as \wib. The codebase is MIT; the
product is New Jersey. Both are deliberate.}

So when \emph{you} (or the agent) work \emph{on} agentty, the rule is: obey the
MIT contract in the code you write, while behaving the New Jersey way in how you
land the change. The seven habits below are exactly that, keyed to real files.

\begin{agentbox}[1 · Ship the good-enough change, then iterate]
Agentty's own system prompt says it: on a vague request (``fix it,'' ``make it
better'') make one reasonable improvement \emph{now} rather than presenting a
menu and waiting. A concrete \texttt{edit} that lands beats an elegant plan that
stalls. \emph{Ship first} at the granularity of a single diff.
\end{agentbox}

\begin{agentbox}[2 · Smallest diff --- \texttt{edit}, not \texttt{write}]
The tool guidance is a \wib\ rule in disguise: use \texttt{edit} (a targeted
substitution, a reviewable diff) over \texttt{write} (dumps the whole file, stalls
the stream). Don't refactor a module when a five-line change will do. Small
changes are cheap to review, revert, and trust.
\end{agentbox}

\begin{agentbox}[3 · Correctness has a floor --- here, the compiler]
\wib\ tolerates \emph{incompleteness}, never \emph{incorrectness of what you
shipped}. In agentty that floor is enforced mechanically: run
\texttt{cmake --build build -j} (the build \emph{is} the test runner). If you
touch a \texttt{static\_assert}-guarded table (\texttt{policy.hpp},
\texttt{spec.hpp}, \texttt{error\_class.hpp}), a wrong edit \emph{won't compile}.
Never route around a proof to ``ship faster.''
\end{agentbox}

\begin{agentbox}[4 · Defer rare cases --- out loud]
Handle what the user will actually hit; \emph{name} the edge-cases you deferred
in your reply. Silent incompleteness is a bug; declared incompleteness is a
roadmap. (This is why agentty surfaces every tool action as its own card\,---\,
nothing important happens invisibly.)
\end{agentbox}

\begin{agentbox}[5 · Measure ``better'' by the user's metric (Waldo)]
Don't chase your own elegance. Agentty already encodes this: retrieval optimises
for \emph{the user's} relevant context, Smart Mode optimises for \emph{their}
cost and latency, not for a maximal answer. Your edit's ``better'' is what serves
the user's goal\,---\,speed, clarity, fitting the repo\,---\,not your scorecard.
\end{agentbox}

\begin{agentbox}[6 · Respect the one style --- reduce switching cost]
Gabriel's lock-in argument, turned inward. \texttt{DESIGN.md}: ``there is one
style.'' So match it exactly\,---\,\texttt{variant} for state, \texttt{expected}
for failure, the dynamism boundary for IO, a \texttt{static\_assert} beside any
new table. A ``more elegant'' pattern that fights the house style is the MIT
mistake: correct in a vacuum, rejected in review.
\end{agentbox}

\begin{agentbox}[7 · Know which hat you're wearing]
The whole trick. \textbf{Writing agentty's code} $\Rightarrow$ \MITs: types,
proofs, consistency, no cut corners. \textbf{Acting as the agent inside it}
$\Rightarrow$ \NJ: smallest diff, ship, iterate. And for genuinely irreversible
work\,---\,auth (\texttt{auth::Credentials}), the permission policy, persistence,
concurrency\,---\,it's \MITs\ \emph{both} ways: get it right before it ships.
\end{agentbox}

\subsection{A quick decision guide}
{\small
\begin{longtable}{@{}>{\RaggedRight}p{6.2cm} >{\RaggedRight}p{7.0cm}@{}}
\toprule
\sffamily\bfseries Situation & \sffamily\bfseries Mode, and why \\
\midrule
\endhead
Vague request, low blast radius (a small feature, a doc, formatting) &
\textcolor{terra}{\bfseries New Jersey.} Ship a reasonable \texttt{edit} now;
iterate. \\
Bug in a common code path &
\textcolor{terra}{\bfseries New Jersey}, compiler as the floor. Smallest correct
diff; \texttt{cmake --build build -j} green before ``done.'' \\
Editing a \texttt{static\_assert}-guarded table (\texttt{policy}, \texttt{spec},
\texttt{error\_class}) &
\textcolor{slate}{\bfseries MIT.} Change the proof \emph{and} the code together;
let the build check you. \\
Auth, credentials, persistence, permissions, concurrency &
\textcolor{slate}{\bfseries MIT / The Right Thing.} Typed, consistent, correct
\emph{before} shipping. No \texttt{EINTR}-style ``deal with it later.'' \\
``Refactor it to be perfect'' &
\textcolor{terra}{\bfseries Push back} toward the smallest useful slice\,---\,but
keep the one style. \\
Edge-case you noticed, user didn't ask &
Handle if cheap; otherwise \textbf{defer and say so}. \\
\bottomrule
\end{longtable}}

\begin{keybox}[The agentty mantra]
\itshape Wear the MIT hat when you \emph{write} the code\,---\,types, proofs, one
style, let the build be the test. Wear the New Jersey hat when you \emph{land} the
change\,---\,smallest \texttt{edit}, measured by the user's yardstick, ship and
iterate, name what you deferred. And for anything irreversible, wear both:
do the Right Thing before it ships.
\end{keybox}

% =============================================================================
\section*{Take-aways in six lines}
\addcontentsline{toc}{section}{Take-aways in six lines}
% =============================================================================
\begin{enumerate}[leftmargin=1.6em]
  \item Simple-and-shipping beats perfect-and-late\,---\,first-mover advantage
        plus lock-in.
  \item The staircase: gain acceptance $\rightarrow$ lower expectations
        $\rightarrow$ grow to almost-right.
  \item ``Worse'' means \emph{simpler and less complete}, never \emph{broken}; a
        ``good-enough'' 50--80\% solution is the target.
  \item The \texttt{EINTR} story is \wib\ made real: shipped simple, spread,
        improved\,---\,but left a permanent wart.
  \item Waldo's rebuttal: it wasn't ``worse'' that won, but \emph{better on the
        metric that mattered}. Measure goodness by the user's ruler.
  \item It's a \emph{force}, not a rule\,---\,Gabriel argued both sides. And
        \texttt{agentty} lives the split: an \MITs\ \emph{codebase} (types,
        compile-time proofs, one style) that ships \NJ\ \emph{behavior} (one lean
        binary, 80\% less context, smallest \texttt{edit}, iterate)\,---\,so wear
        the right hat for the task, and both hats for anything irreversible.
\end{enumerate}

\vfill
\begin{center}
{\color{brass}\rule{0.35\linewidth}{0.4pt}}\\[6pt]
{\sffamily\scriptsize\color{faint}
Sources: R.\,P.\ Gabriel, \emph{The Rise of ``Worse is Better''} (1989) $\cdot$
Jim Waldo, \emph{``Worse is Worse''} (Artima, 2003) $\cdot$
J.\ Haberman, \emph{``EINTR and PC loser-ing''} (2011) $\cdot$
Gabriel, \emph{``Worse Is Better Is Worse''} \& \emph{``Is Worse Really Better?''} (2000)\\
and the Wikipedia article ``Worse is better'' and its further reading
(Gancarz, Graham, Mayer). The \texttt{agentty} section draws on the project's own
\texttt{docs/DESIGN.md}, \texttt{ARCHITECTURE.md}, \texttt{WHY-NOT-RUST.md}, and
\texttt{README.md}.}
\end{center}

\end{document}
